\section{Data and Simulated Samples}%
\label{ch:data-sim}

%------------------------------------------------------------------------------

The analysis presented in this note uses the full dataset collected by the ATLAS
detector during Run2 operations (2015-2018). This consists of data from proton-proton  %chktex 8
collision at a centre-of-mass energy of \qty{13}{\TeV} and has a total integrated
luminosity of \lumi.

%-------------------------------------------------------------------------------
\subsection{Data sample}%
\label{sec:datamc:data}
%-------------------------------------------------------------------------------


Only data collected during stable beam LHC operations and with the ATLAS detector
fully functioning is included. The partial integrated luminosities for each year
of operation\footnote{\url{https://twiki.cern.ch/twiki/bin/view/Atlas/LuminosityForPhysics}},
as well as the corresponding Good Run Lists (GRLs) used are reported in \cref{tab:datamc:lumi}:

\begin{table}[htbp]
	\centering
  \caption{Integrated luminosity per year.}%
	\label{tab:datamc:lumi}
	\begin{tabular}{cSl}
		\toprule
		Year & {Int.\ lumi.\ (\si{\per\pb})} & GRL \\
		\midrule
		2015 & 3244.54   & data15\_13TeV/20170619/physics\_25ns\_21.0.19.xml \\
		2016 & 33402.2	 & data16\_13TeV/20180129/physics\_25ns\_21.0.19.xml \\
		2017 & 44630.6	 & data17\_13TeV/20180619/physics\_25ns\_Triggerno17e33prim.xml \\
		2018 & 58791.6   & data18\_13TeV/20190318/physics\_25ns\_Triggerno17e33prim.xml  \\
		\midrule
		Run 2 & 140068.94 & \\
		\bottomrule
	\end{tabular}
\end{table}

%-------------------------------------------------------------------------------
\subsection{Monte Carlo simulated samples}%
\label{sec:datamc:mc} 
%-------------------------------------------------------------------------------  

The contribution of SM processes is evaluated from Monte Carlo (MC) simulated
samples~\cite{SOFT-2010-01}.
The corresponding dataset IDs are given in~\cref{app:samples}.

The effect of multiple interactions in the same and neighbouring bunch
crossings (\pileup) was modelled by overlaying the simulated hard-scattering event with
inelastic proton--proton (\(pp\)) events generated with \PYTHIA[8.186]~\cite{Sjostrand:2007gs}  %chktex 8
using the \NNPDF[2.3lo] set of parton distribution functions (PDF)~\cite{Ball:2012cx} and the
A3 set of tuned parameters~\cite{ATL-PHYS-PUB-2016-017}.

The detector response is simulated with \texttt{Geant4}~\cite{Agostinelli:2002hh}.
As a full simulation of the detector response is computationally expensive for
some of the estimates of modelling uncertainties, the fast-simulation package
\textsc{AtlFast-II}~\cite{SOFT-2010-01} is used, which parametrises hadronic
showers in the electromagnetic and hadronic calorimeter to speed up the simulation.

The Monte Carlo (MC) events were weighted to reproduce the
distribution of the average number of interactions per bunch crossing
(\(\left<\mu \right>\)) observed in the data. The \(\left<\mu \right>\)
value in data was rescaled by a factor of \(1.03\pm 0.04\) to improve
agreement between data and simulation in the visible inelastic
proton--proton (\(pp\)) cross-section~\cite{STDM-2015-05}.  %chktex 8

The signal samples are defined as the sum of the \ttbar and the \tW samples,
simulated with the same combination of hard process generators and parton showers.
The description of the MC samples is taken from the official PMG documentation.

Only prompt leptons are taken from Monte Carlo. They originate from decays of
\Zboson, \Wboson, and \Hboson bosons, or from prompt
leptonic \( \tau \) decays. Fake lepton contributions come from processes such
as in-flight decays of mesons, jets reconstructed as leptons, and electron
photon conversions. More details on fake estimation are provided in \cref{sec:bkg-estimation}.

%-------------------------------------------------------------------------------
\subsubsection{Nominal signal sample}%
\label{sec:datamc:mc_nominal_signal} 
%-------------------------------------------------------------------------------  

\paragraph*{\ttbar production}

The production of \ttbar events was modelled using the
\POWHEGBOX[v2]~\cite{Frixione:2007nw,Nason:2004rx,Frixione:2007vw,Alioli:2010xd}
generator, which provided matrix elements at next-to-leading
order~(NLO) in the strong coupling constant \alphas, and used the
\NNPDF[3.0nlo]~\cite{Ball:2014uwa} set of parton distribution functions~(PDFs) as input.
The \hdamp parameter, which influences the matching of \POWHEG hard scatter events
and parton shower, and effectively regulates the high-\pt radiation against which the
\ttbar system recoils, was set to 1.5\,\mtop~\cite{ATL-PHYS-PUB-2016-020}.
The functional form of the renormalisation and factorisation scales was
set to the default scale \(\sqrt{\mtop^{2} + p_\mathrm{T\ top}^2}\).
The events were interfaced with
\PYTHIA[8.230]~\cite{Sjostrand:2014zea} for the parton shower and
hadronisation, using the A14 set of tuned
parameters~\cite{ATL-PHYS-PUB-2014-021} and the \NNPDF[2.3lo]
set of PDFs~\cite{Ball:2012cx}.
The decays of bottom and charm hadrons were simulated using the
\EVTGEN[1.6.0] program~\cite{Lange:2001uf}.

The \ttbar sample was normalised to the cross-section prediction at
next-to-next-to-leading order (NNLO) in QCD including the resummation of
next-to-next-to-leading logarithmic (NNLL) soft-gluon terms calculated using
\TOPpp[2.0]~\cite{Beneke:2011mq,Cacciari:2011hy,Baernreuther:2012ws,Czakon:2012zr,Czakon:2012pz,Czakon:2013goa,Czakon:2011xx}.
For proton--proton collisions at a centre-of-mass energy of  % chktex 8
\(\rts = \qty{13}{\TeV}\), this cross-section corresponds to
\(\sigma{(\ttbar)}_{\mathrm{NNLO+NNLL}} = 832^{+41}_{-51}\,\unit{\pb}\) using
a top-quark mass of \(\mtop = \qty{172.5}{\GeV}\).
The uncertainties in the cross-section due to the PDF and \alphas were
calculated using the \PDFforLHC[15] prescription~\cite{Butterworth:2015oua}
with the \MMHT[nnlo]~\cite{Harland-Lang:2014zoa},
\CT[14nnlo]~\cite{Dulat:2015mca} and \NNPDF[3.0lo]~\cite{Ball:2014uwa}
PDF sets in the five-flavour scheme.

The uncertainty due to initial-state radiation (ISR) was estimated by
comparing the nominal \ttbar sample with two additional
variations. The \hdamp value was increased to 3.0\,\mtop and compared to the
nominal value. Separately the Var3c parameter of the A14 parton shower tune was varied,
largely corresponding to the variation of \alphas for ISR~\cite{ATL-PHYS-PUB-2014-021}.
The impact of final-state radiation (FSR) was evaluated by
varying the renormalisation scale for emissions from the
parton shower up and down by a factor of two.

The statistically enhanced \ttbar sample, using slices in particle-level 
 \MET (\( \MET > \qty{200}{\GeV} \)) and \HT
(\( \HT > \qty{600}{\GeV} \& \MET < \qty{200}{\GeV} \)), is used in addition
to the nominal sample (\(\MET < \qty{200}{\GeV} \& \HT < \qty{600}{\GeV}\)).
Orthogonal cuts between the samples ensure no overlap.
The samples are listed in \cref{app:samples}. When an additional sample is needed
for systematics variation it is not statistically enhanced.

Additionally, all \ttbar and \tW events with at least one lepton in the final state (including dilepton events that pass the selection due to out-of-acceptance effects) are taken into account in the signal
region due to the use of fully inclusive \ttbar samples.

Two additional \ttbar variations are used in the final
comparison. One which is obtained by reweighting the nominal \ttbar sample
to a NNLO QCD + NLO EW prediction~\cite{Czakon:2017wor} at parton-level,
using the setting proposed in the reference. The reweighting is performed
on the three kinematic variables \(p_T(t)\), \(m(t\bar{t})\) and \(p_T(t\bar{t})\)
of the top quarks in the MC samples after initial- and final-state radiation.
The predictions for \(p_T(t)\) and \(m(t\bar{t})\) are calculated at NNLO in
QCD with NLO electroweak (EW) corrections with the NNPDF3.0qed PDF set using
the dynamic renormalisation and factorisation scales \(m_T(t)/2\) for
\(p_T(t)\) and \(H_T/4\) for \(m(t\bar{t})\) as proposed in Ref.~\cite{Czakon:2017wor}.
The prediction for \(p_T(t\bar{t})\) is calculated at NNLO in QCD with the
NNPDF3.0 PDF set and with \(\mu_r\) and \(\mu_f\) set to \(H_T/4\).
The reweighting was performed iteratively, such that at the end of the procedure
the reweighed MC sample agrees well with the higher-order prediction for each
of the three variables. The re-weighted predictions themselves are not
equivalent to complete NNLO+parton shower calculations and are only indicative
of the NNLO contributions on the measured observables. The short name for this
sample is marked as  (NNLO rew.).
The second alternative MC sample  to model \ttbar processes at NNLO QCD
incorporates all-order radiative corrections through parton-shower simulations
as detailed in Ref.~\cite{Mazzitelli:2020jio}.
The sample is normalised to the same NNLO+NNLL in pQCD calculated with
the Top++ 2.0 program, as previously mentioned.
Two  different
configurations for the \(p_T,hard\) and \(p_T,def\) parameters were generated:
one with \(p_T,hard = 0\) and \(p_T,def = 2\), and another with \(p_T,hard = 1\)
and \(p_T,def = 1\). In all subsequent figures and tables, predictions based
on these samples are referred to as `MiNNLO' and `MiNNLO (alt)', respectively.





\paragraph*{\tW production}

Single-top \(tW\) associated production was modelled using the
\POWHEGBOX[v2]~\cite{Re:2010bp,Nason:2004rx,Frixione:2007vw,Alioli:2010xd}
generator, which provided matrix elements at next-to-leading
order~(NLO) in the strong coupling constant \alphas\ and made use of the five-flavour
scheme with the \NNPDF[3.0nlo]~\cite{Ball:2014uwa} parton
distribution function~(PDF) set.
The functional form of the renormalisation and factorisation scales was
set to the dynamic scale \( \HT / 2 \)~\cite{ATL-PHYS-PUB-2021-042}.
The diagram removal scheme~\cite{Frixione:2008yi} was employed to handle the
interference with \ttbar production~\cite{ATL-PHYS-PUB-2016-020}.
The events were interfaced with \PYTHIA[8.230]~\cite{Sjostrand:2014zea}
using the A14 tune~\cite{ATL-PHYS-PUB-2014-021} and the \NNPDF[2.3lo] PDF set.
The decays of bottom and charm hadrons were simulated using the
\EVTGEN[1.6.0] program~\cite{Lange:2001uf}.

The inclusive cross-section was corrected to the theory prediction
calculated at NLO in QCD with NNLL soft-gluon
corrections~\cite{Kidonakis:2021vob}.  For proton--proton  % chktex 8
collisions at a centre-of-mass energy of \(\rts = \qty{13}{\TeV}\), this
cross-section corresponds to \(\sigma{(tW)}_{\mathrm{NLO+NNLL}} = 79.3 \pm 2.9\,\unit{\pb}\),
using a top-quark mass of \(\mtop = \qty{172.5}{\GeV}\).
The uncertainties in the cross-section due to the PDF and \alphas were
calculated using the \PDFforLHC[15] prescription~\cite{Butterworth:2015oua}
with the \MMHT[nnlo]~\cite{Harland-Lang:2014zoa},
\CT[14nnlo]~\cite{Dulat:2015mca} and \NNPDF[3.0lo]~\cite{Ball:2014uwa}
PDF sets in the five-flavour scheme.

Like in the \ttbar sample uncertainty due to initial-state radiation (ISR)
was estimated by separately varying the \hdamp parameter and
the Var3c parameter of the A14 tune
as described in Ref.~\cite{ATL-PHYS-PUB-2017-007}. The impact of
final-state radiation (FSR) was evaluated by varying the renormalisation scale
for emissions from the parton shower up or down by a factor two.

This sample is displayed on plots and in tables as PH+PY8 (DR) or \tW (DR).

%-------------------------------------------------------------------------------  
\subsubsection{Alternative signal samples}%
\label{sec:datamc:mc_alt_signal} 
%-------------------------------------------------------------------------------  

Several additional \ttbar and \tW samples are considered for evaluating modelling systematics and for comparison with data. 

\begin{itemize}
	\item \tW PH+PY8 (DS): similar to the nominal PH+PY8 sample using dynamic scales but using
    the diagram subtraction scheme~\cite{Frixione:2008yi,ATL-PHYS-PUB-2016-020}.
    Combined with the nominal \ttbar sample, it is presented in legends and tables as PH+PY8 (DS).
	\item \tW PH+PY8 with fixed scales (both DR and DS):
    similar to the nominal PH+PY8 sample but with the renormalisation and
    factorisation scales fixed to the default scale, \(\mtop = \qty{172.5}{\GeV}\).
    Combined with the nominal \ttbar sample, they are presented in legends
    and tables as PH+PY8 (DRF) and PH+PY8 (DSF) respectively. See Ref.~\cite{ATL-PHYS-PUB-2021-042} for detailed on improvements of using dynamic scales.
	\item \ttbar and \tW PH+PY8 (\(\hdamp = 3 \mtop \)):
    produced with the same settings as the nominal PH+PY8 sample but with the 
	\hdamp parameter set to \( 3 \mtop \). The short name used in the paper is PH+PY8 (\hdamp).
	\item \ttbar and \tW PH+H7: events were generated using the same \POWHEGBOX[v2] set-up
    using the \NNPDF[3.0nlo] PDF, interfaced to \HERWIG[7.1.3]~\cite{Bellm:2015jjp,Bahr:2008pv},
    which includes an angular-ordered parton-shower model with a different hadronisation model. This sample was generated using the \HERWIG[7.1]
    default set of tuned parameters and the \MMHT[lo] PDF set~\cite{Harland-Lang:2014zoa}.
    For the \tW process, events were generated with the \POWHEGBOX[v2] generator were
    interfaced to \HERWIG[7.2].
    The decays of the bottom and charm hadrons were simulated using the \EVTGEN[1.6.0]
    program~\cite{Lange:2001uf}. The short name used in the paper is PH+H7.
	\item \ttbar and \tW PH+PY8 (\pthard): similar to the nominal PH+PY8 sample
    but using \( \pthard = 1 \) instead of  \( \pthard = 0 \)~\cite{10.21468/SciPostPhys.12.1.010}.
    This parameter regulates the definition of the vetoed region of the showering,
    important to avoid hole/overlap in the phase space filled by \POWHEG and \PYTHIA.
    The short name for this sample is PH+PY8 (\pthard).
	\item \ttbar and \tW PH+PY8 with varied top-quark mass:
    mass variations of \( \pm 0.5 \mtop \) for both \ttbar and \tW are used.
	\item \ttbar PH+PY8+MadSpin: similar to the nominal PH+PY8 sample but with
    \madspin~\cite{Artoisenet:2012st} program used to decay the top quarks.
    The decays of the bottom and charm hadrons were simulated using the \EVTGEN[1.7.0]
    program. This sample is combined with the nominal \tW sample and is named PH+PY8+MadSpin.
    \item \ttbar with recoil to top-quark: similar to the nominal PH+PY8 sample
    but the modelling of the parton shower employs the recoiling against top quarks instead of \bquarks.
    \item \ttbar and \tW aMC@NLO+PY8: \MGNLO generator is used instead of \PYTHIA for the matrix element calculation.
\end{itemize}

All the alternative samples are normalised to the same total cross-section
as the nominal, except for the top-quark mass variations, which are normalised
to the corresponding higher-order cross-sections.


%-------------------------------------------------------------------------------
\subsubsection{Background samples}%
\label{sec:datamc:mc_background}
%------------------------------------------------------------------------------- 

Below is the list of the background samples used in the analysis. The dominant
background are the \Wjets events. Other backgrounds are grouped in categories.
In the plots and tables ``\OtherT'' consists of \(s\)-channel and \(t\)
channel single-top production and the \ttV / \ttH samples.
``\OtherV''  consists of \Zjets and diboson samples.
Treatment of misidentified and fake leptons, denoted also as ``Mis-ID lep.''  % chktex 38
is described in more detail in \cref{sec:bkg-estimation}.

\paragraph*{\Wjets and \Zjets production}

The production of \(V+\)jets was simulated with the \SHERPA[2.2.11]~\cite{Bothmann:2019yzt}
generator. In this set-up, NLO-accurate matrix elements for up to two additional partons,
and LO-accurate matrix elements for up
to five additional partons were calculated with the Comix~\cite{Gleisberg:2008fv} and
\OPENLOOPS~\cite{Buccioni:2019sur,Cascioli:2011va,Denner:2016kdg} libraries.
The default \SHERPA parton showering algorithm~\cite{Schumann:2007mg} based on
the Catani--Seymour dipole factorisation and the cluster hadronisation model~\cite{Winter:2003tt}  % chktex 8
was used. The dedicated set of tuned parameters developed by the
\SHERPA authors and the Hessian \NNPDF[3.0nnlo] PDF set~\cite{Ball:2014uwa,Carrazza:2015aoa} were employed.

The NLO matrix elements for a given jet multiplicity were matched to the parton
shower (PS) using a colour-exact variant of the MC@NLO
algorithm~\cite{Hoeche:2011fd}. Different jet multiplicities were then merged
into an inclusive sample using an improved CKKW matching
procedure~\cite{Catani:2001cc,Hoeche:2009rj} which was extended to NLO
accuracy using the \MEPSatNLO prescription~\cite{Hoeche:2012yf}. The merging threshold
was set to \qty{20}{\GeV}.

Electroweak virtual corrections are included in the 2.2.11 configuration,
and are available as an alternative generator weight based on the
NLO \(EW_\mathrm{virt}\) approach using the exponentiated
approach~\cite{Kallweit:2014xda,Kallweit:2015dum}.

Uncertainties from missing higher orders were
evaluated~\cite{Bothmann:2016nao} using seven variations of the QCD
factorisation and renormalisation scales in the matrix elements, changing them by
factors of \(0.5\) and \(2\), avoiding variations in opposite directions.
An envelope is taken of all variations.

Uncertainties in the nominal PDF set were evaluated using 100 replica
variations. Additionally, the results were cross-checked using the
central values of the \CT[14nnlo]~\cite{Dulat:2015mca} and
\MMHT[nnlo]~\cite{Harland-Lang:2014zoa} PDF sets. The effect of the uncertainty
in the strong coupling constant \( \alphas \) was assessed by variations
of \(\pm 0.001\).

The \(V\)+jets samples were normalised to a next-to-next-to-leading-order (NNLO)
prediction~\cite{Anastasiou:2003ds}.


\paragraph*{Diboson production}

Samples of diboson final states (\(VV\)) were simulated with the
\SHERPA[2.2.1] or 2.2.2~\cite{Bothmann:2019yzt} generator depending on the process,
%~\footnote{This is an admixture of 2.2.1 and above versions, so the version should be kept generic to avoid confusions. As an
%alternative, the sentence can be modified indicating that samples are simulated with the \SHERPA[2.2.1] or 2.2.2 depending on the process.}
including off-shell effects and Higgs boson contributions, where appropriate.
Fully leptonic final states and semileptonic final states, where one boson
decays leptonically and the other hadronically, were generated using
matrix elements at NLO accuracy in QCD for up to one additional parton
and at LO accuracy for up to three additional parton
emissions. Samples for the loop-induced processes \(gg \to VV\) were
generated using LO-accurate matrix elements with up to one
additional parton emission for both the cases of fully leptonic and
semileptonic final states. The matrix element calculations were matched
and merged with the \SHERPA parton shower based on the Catani--Seymour  % chktex 8
dipole factorisation~\cite{Gleisberg:2008fv,Schumann:2007mg} using the MEPS@NLO
prescription~\cite{Hoeche:2011fd,Hoeche:2012yf,Catani:2001cc,Hoeche:2009rj}.
The virtual QCD corrections were provided by the
\OPENLOOPS library~\cite{Buccioni:2019sur,Cascioli:2011va,Denner:2016kdg}. The
\NNPDF[3.0nnlo] set of PDFs was used~\cite{Ball:2014uwa}, along with the
dedicated set of tuned parton-shower parameters developed by the
\SHERPA authors.

No overlap with the \Wjets and \Zjets samples is present, as hadronic decays
are prompt vector boson decays and not associated jet production. Hence the kinematics
of the corresponding events are different.

\paragraph*{Other single top production}

Single-top \(t\)-channel and \(s\)-channel productions were modelled using the
\POWHEGBOX[v2]~\cite{Frederix:2012dh,Alioli:2009je,Nason:2004rx,Frixione:2007vw,Alioli:2010xd}
generator at NLO in QCD using the four-flavour scheme for \(t\)- and five-flavour scheme for \(s\)-channel and the
corresponding \NNPDF[3.0nlo] set of PDFs~\cite{Ball:2014uwa}.  The events were
interfaced with \PYTHIA[8.230]~\cite{Sjostrand:2014zea} using the A14
tune~\cite{ATL-PHYS-PUB-2014-021} and the \NNPDF[2.3lo] set of
PDFs~\cite{Ball:2012cx}.


\paragraph*{\ttV production}

The production of \ttV events was modelled using the \MGNLO[2.3.3]~\cite{Alwall:2014hca}
generator, which provided matrix elements at
next-to-leading order~(NLO) in the strong coupling constant \alphas
and used the \NNPDF[3.0nlo]~\cite{Ball:2014uwa} set of parton distribution
functions~(PDFs). The functional form of the renormalisation and
factorisation scales was set to the default of
\(0.5 \times \sum_i \sqrt{m^2_i+p^2_{\mathrm{T},i}}\),
where the sum runs over all the particles
generated in the matrix element calculation. Top quarks were decayed
at LO using \MADSPIN~\cite{Frixione:2007zp,Artoisenet:2012st} to
preserve spin correlations. The events were interfaced with
\PYTHIA[8.210]~\cite{Sjostrand:2014zea} for the parton shower and
hadronisation, using the A14 set of tuned
parameters~\cite{ATL-PHYS-PUB-2014-021} and the
\NNPDF[2.3lo]~\cite{Ball:2014uwa} PDF set.
The decays of bottom and charm hadrons were simulated using the \EVTGEN[1.2.0] program~\cite{Lange:2001uf}.


\subsubsection*{\ttH production}

The production of \ttH events was modelled using the
\POWHEGBOX[v2]~\cite{Frixione:2007nw,Nason:2004rx,Frixione:2007vw,Alioli:2010xd,Hartanto:2015uka}
generator at NLO with the \NNPDF[3.0nlo]~\cite{Ball:2014uwa} PDF set.
The events were interfaced to \PYTHIA[8.230]~\cite{Sjostrand:2014zea}~using
the A14 tune~\cite{ATL-PHYS-PUB-2014-021} and the
\NNPDF[2.3lo]~\cite{Ball:2014uwa} PDF set. The decays of bottom and charm hadrons
were performed by \EVTGEN[1.6.0]~\cite{Lange:2001uf}.
